\documentclass[11pt]{article}
\usepackage[margin=1in]{geometry}
\usepackage{booktabs}
\usepackage{amsmath}
\usepackage{hyperref}
\usepackage{longtable}
\title{Independent Replication of Matsuya et al.\ (2018): The Integrated
Microdosimetric-Kinetic (IMK) Model of DNA-Targeted and Non-Targeted Effects}
\author{LUCID Replication Project (rick-stevens-ai)}
\date{2026-05-28 (report), backfill 2026-07-06}
\begin{document}
\maketitle

\section{Paper summary}
Matsuya et al.\ (2018, \emph{Sci.\ Rep.}\ 8:4849; CC-BY;
DOI:10.1038/s41598-018-23202-y) present the \textbf{Integrated
Microdosimetric-Kinetic (IMK) model}, a mechanistic ODE-based framework
that couples classical DNA-targeted effects (DTE) with non-targeted effects
(NTE, i.e.\ bystander signalling). The model has three coupled layers:
(i) fast Ca$^{2+}$ transient + slow NO signal kinetics (Eq.~9); (ii) direct
DSB production with two-rate biphasic repair (Eqs.~2--3, 10--12), where
NTE-induced ``persistent long-lived'' (PLL) lesions repair
$\sim\!65\times$ slower than DTE DSBs; and (iii) LQ-form survival with an
NTE contribution that saturates through a Poisson-occluded ``hit fraction''
$f_h(D)=1-\exp(-N_h)$ (Eqs.~7--8, 15). The model is fitted against digitised
data from V79-379A, T-47D, HPV-G, E48, CHO-K1 and MRC-5. Central biological
claims: (a) low-dose hyper-radiosensitivity (HRS) emerges naturally when
non-hit-cell repair rate $c_b$ is lower than hit-cell $c$; (b) MTBE
(medium-transfer bystander effect) curves fit Eq.~25 in HPV-G/E48;
(c) a factor-of-4 span of $c_b$ produces monotonically deeper HRS
(Fig.~5B). The paper's most striking prediction is that reduced repair
efficiency in \emph{non-hit} cells is \emph{necessary} for HRS.

\section{Replication scope}
We independently reimplemented Eqs.~1--26 and SI-1--SI-10 in pure Python
(NumPy + SciPy, CPU-only, no GPU). We used hand-digitised data from the
paper's own Figs.~2--4 plus the printed parameters in Tables 1--2 as our
comparison target. \textbf{No raw experimental data} from the cited
underlying works (Lyng 2002, Han 2007, Marples 1993--99, Edin 2007--12,
Mothersill 2000, Liu 2006/2007, Ojima 2011, Chalmers 2004) was accessed.
No public IMK source was located (checked github.com/topas-nbio and
Y. Matsuya's software footprint as of 2026-05-28).

\section{Method}
\begin{itemize}
    \item PDF: paper.pdf sha256
      \texttt{eda2685eb056e546a68bfd2f937db32d79ee4796775cbd124078bb0b515659d8}.
    \item Code: \texttt{code/imk\_model.py}, \texttt{code/make\_figures.py},
      \texttt{code/reference\_data.py}.
    \item Numerics: closed-form for Eqs.~7--9, 15, 25; analytic
      solutions for the biphasic DSB kinetics; bounded nonlinear
      least-squares (SciPy \texttt{trf}) on log-SF for the
      independent parameter refit. No MCMC.
    \item Runtime: $\sim\!4$~s CPU (numpy 1.x + scipy 1.x + matplotlib).
\end{itemize}

\section{Claim-by-claim results (vs.\ paper)}
Ten explicit numeric claims were extracted from the paper and tested.
Status keys: REPLICATED = matched within $\sim$10\% on the digitised data;
SPOT-CHECK = curve shape matched, tight metric limited by digitisation;
PARTIAL = matched on part of the data; CONTRADICTED = different answer.
Full narrative table with equation references is in \texttt{../REPORT.md}.
Summary:

\begin{longtable}{p{0.06\linewidth} p{0.30\linewidth} p{0.14\linewidth} p{0.32\linewidth} p{0.10\linewidth}}
\toprule
\# & Claim (compact) & Eq. & Our finding & Status \\
\midrule
1 & NTE emission $\sim$ LQ with $f_h(D)$ & 7--8 & LQ shape exact by construction; $1/\sqrt{\beta_b}$ text vs.\ Table 2 internally inconsistent for V79 & PARTIAL \\
2 & Ca$^{2+}$ + NO kinetics (Eq.~9) & 9 & Analytic Eq.~9 correct; sparse digitisation gives negative R$^2$ on Ca but shape matches Lyng 2002 & SPOT-CHECK \\
3 & NTE repair $\sim$65$\times$ slower than TE (MRC-5) & 2,3,10,12 & Reproduced exactly (64.6$\times$) & \textbf{REPLICATED} \\
4 & V79-379A HRS shoulder near 0.2--0.5~Gy & 4,17,19 & HRS dip present in $-\ln S/D$; digitisation-limited R$^2$ & SPOT-CHECK \\
5 & T-47D survival with HRS & 4,17,19 & HRS shape OK; absolute $S(2\text{ Gy})\approx 0.67$ vs.\ published $\sim$0.1 with printed $\beta_0=0.029$ & PARTIAL \\
6 & HPV-G, E48 MTBE curves (Eq.~25) & 25--26 & E48 R$^2$=0.97; HPV-G plateau $\exp(-\delta_m)\approx 0.41$ matches curve shape not coarse points & PARTIAL \\
7 & PARP-inhibited CHO-K1 deepens HRS (Fig.~4) & 19 & Sham vs.\ inhibited separation reproduced qualitatively & \textbf{REPLICATED (qual.)} \\
8 & Derived $c_b=0.155$~h$^{-1}$ in CHO-K1 & post-26 & Direct back-solve gives 0.155 & \textbf{REPLICATED} \\
9 & $c_b$ scan ($\times$4,1,$\tfrac{1}{2}$,$\tfrac{1}{4}$) deepens HRS (Fig.~5B) & 16,19 & Reproduced (\texttt{fig6\_hrs\_repair\_scan.png}) & \textbf{REPLICATED} \\
10 & Max LLs/nucleus $\approx$ 0.23 for T-47D (Fig.~5A) & 15 & Eq.~15 max is $\delta/4$; T-47D $\delta=0.172 \Rightarrow$ max 0.043 vs.\ 0.23 & \textbf{CONTRADICTED} \\
\bottomrule
\end{longtable}

\noindent \textbf{Coverage}: 10/10 explicit numeric claims attempted.
\textbf{Agreement}: 4 REPLICATED (incl.\ 1 qualitative), 2 SPOT-CHECK,
3 PARTIAL, 1 CONTRADICTED $\Rightarrow \approx 70\%$.

\section{What worked}
\begin{itemize}
  \item Every closed-form equation (7, 8, 9, 15, 25) reproduces exactly:
    the algebra of the IMK model is unambiguous and Eq.~9 is standard
    exponential-difference kinetics.
  \item The signature qualitative predictions --- HRS shoulder,
    $c_b$-scan monotonic deepening of HRS, PARP-inhibition-vs-sham
    ordering, and NTE:TE repair ratio $\sim$65$\times$ --- all
    reproduced without parameter tuning.
  \item The algebraically derived $c_b=0.155$~h$^{-1}$ from the CHO-K1
    fit constants matches the paper's quoted value.
\end{itemize}

\section{What didn't}
\begin{itemize}
  \item \textbf{Claim 1 (V79 $1/\sqrt{\beta_b}$).} Table 2 gives
    $\beta_b=0.396$ Gy$^{-2}$, so $1/\sqrt{\beta_b}=1.59$~Gy, but the
    text after Eq.~7 states 5.03~Gy. This is a
    \emph{paper-internal inconsistency}: either the exponent in
    Table 2 is off by $10\times$ ($\beta_b\approx 0.0396$) or the text
    number is a typo. Our arithmetic is correct either way.
  \item \textbf{Claim 5 (T-47D absolute SF).} With the printed
    $\alpha_0=0.129, \beta_0=0.029$, our model predicts
    $S(2\text{ Gy})\approx 0.67$; the published Fig.~2(D) shows
    $\sim$0.1. Either $\beta_0$ is a typo ($\sim$0.06 would close the
    gap) or the plotted curve uses additional NTE terms that Eq.~15
    should not admit at $\delta=0.17$.
  \item \textbf{Claim 6 (HPV-G MTBE).} The published parameters
    ($\delta_m=0.902$, $\alpha_b=30.9$) drive Eq.~25 to a plateau
    $S\to\exp(-\delta_m)\approx 0.41$ at very low dose. Curve
    \emph{shape} matches, but R$^2$ on our coarse digitisation is 0.32.
  \item \textbf{Claim 10 (T-47D LL max).} Eq.~15 has a strict
    theoretical maximum of $\delta/4$. To reach 0.23 (Fig.~5A) requires
    $\delta\approx 0.92$, which is HPV-G's value. Most parsimonious
    explanation: legend swap between V79/HPV-G/T-47D in Fig.~5A.
  \item \textbf{Independent V79 fit lands in a different parameter
    basin} from the paper (R$^2>0.999$ on log-SF but different
    $\alpha_0,\beta_0,\alpha_b,\beta_b,\delta$). This confirms known
    parameter degeneracy in the 5--7-parameter IMK model.
\end{itemize}

\section{Critique}
\begin{enumerate}
  \item \textbf{Digitised-only comparison.} Every metric here compares
    against digitised paper figures, not against raw underlying
    experiments (Lyng, Han, Marples, Ojima, etc.). A truly independent
    replication would refit IMK to the primary microscopy/clonogenic
    data. Our ``REPLICATED'' verdict is therefore explicitly
    ``the paper's math is internally consistent with its own reported
    parameters,'' not ``the paper's biology is confirmed against the
    literature it cites.''
  \item \textbf{Parameter degeneracy is a first-order problem, not an
    aside.} A 5--7-parameter model that has multiple basins with
    R$^2 > 0.999$ on log-SF cannot support strong quantitative
    biological claims (like ``$c_b$ must be 0.155~h$^{-1}$'').
    The paper does not quote credible intervals or a covariance
    matrix; neither do we.
  \item \textbf{No MCMC.} The paper cites $10^7$ MC samples for
    parameter estimation; we used bounded NLS. Neither method
    provides posterior uncertainty in the form a modern
    identifiability analysis would require (profile likelihoods,
    Sobol indices, or Bayesian credible intervals).
  \item \textbf{Bystander-signal dynamics not independently tested.}
    We reproduced Eq.~9's kinetics against digitised Ca$^{2+}$ and
    NO traces but did \emph{not} simulate a spatial-temporal cell
    population or model cell-cell coupling. The IMK model in this
    replication is well-mixed and single-population; the
    ``non-targeted effect'' is folded into $f_h(D)$ and $\delta$
    rather than emerging from paracrine signalling.
  \item \textbf{Low-dose HRS reproduction is model-only.} We
    reproduced the model prediction that $c_b<c$ produces HRS.
    We did \emph{not} demonstrate that the HRS observed in wet-lab
    V79 or T-47D data \emph{requires} NTE-mediated slow repair, as
    opposed to induced-repair (IR) models, PLD-repair variants, or
    simple LQ+threshold. This is the paper's central biological
    claim, and it is untestable against digitised curves alone.
  \item \textbf{Paper self-inconsistencies (Claims 1, 5, 10) look
    like typos, not fundamentals}, but they were not caught by peer
    review; a real replication ecosystem would produce corrigenda,
    not just replication reports.
  \item \textbf{Scope narrow.} We covered clonogenic survival, DSB
    kinetics, and MTBE. We did not touch chromosome aberrations,
    micronuclei, apoptosis, or transformation --- all of which the
    IMK modelling literature claims to address.
\end{enumerate}

\section{Open Questions}
See \texttt{report/open\_questions.json} for machine-readable form.
Five questions of genuine radiobiological interest:
\begin{enumerate}
  \item Is $c_b<c$ really \emph{necessary} for HRS, or is IMK just
    one of several models that can fit it?
  \item Does the paper's LQ-form NTE emission (Eq.~7) survive
    replacement by mechanistic paracrine signalling ODEs?
  \item Are Table~2's cell-line parameters identifiable, or are they
    a single mode of a multi-modal posterior?
  \item Does the 65$\times$ NTE:TE repair ratio hold across cell
    lines beyond MRC-5, or is it an MRC-5 idiosyncrasy?
  \item Does the model's HRS prediction generalise to high-LET
    radiation (protons, $\alpha$, C-ion), where the Poisson-hit
    fraction $f_h$ saturates much faster?
\end{enumerate}

\section{Reproduction}
\begin{verbatim}
cd code/
python3 imk_model.py       # smoke test
python3 make_figures.py    # figures/*.png + results/summary.json
\end{verbatim}

\section{Verdict}
\textbf{REPLICATED} for the IMK model's mathematical structure and its
qualitative biology (HRS shoulder, $c_b$-scan monotonicity, PARP
ordering, NTE:TE repair ratio, MTBE curve shape). \textbf{PARTIAL} on
absolute survival curves where paper-internal inconsistencies (Claims 1,
5, 10) prevent a clean numerical match. \textbf{No falsification} of the
paper's central hypothesis. However, digitised-only comparison and 5--7
degenerate parameters mean this replication supports the paper's
\emph{model} more than its \emph{biology}.

\end{document}
